\section*{Hoofdstuk \ref{ch:g7:fractions} --- Breuken: vergelijken en optellen}

\begin{solution}{exo:g7:fractions:1}
$\dfrac{2}{5} = \dfrac{8}{20}$; \quad
$\dfrac{18}{24} = \dfrac{3}{4}$; \quad
$\dfrac{7}{3} = \dfrac{28}{12}$; \quad
$\dfrac{5}{9} = \dfrac{20}{36}$.
\end{solution}

\begin{solution}{exo:g7:fractions:2}
$\dfrac{12}{16} = \dfrac34$; \quad
$\dfrac{30}{45} = \dfrac23$; \quad
$\dfrac{27}{36} = \dfrac34$; \quad
$\dfrac{48}{12} = 4$.
\end{solution}

\begin{solution}{exo:g7:fractions:3}
$\dfrac{15}{35} = \dfrac37$ en $\dfrac{21}{49} = \dfrac37$: gelijk.

$\dfrac{16}{28} = \dfrac47$ en $\dfrac{20}{36} = \dfrac59$: om te vergelijken,
$\frac47 = \frac{36}{63}$ en $\frac59 = \frac{35}{63}$ --- niet gelijk.
\end{solution}

\begin{solution}{exo:g7:fractions:4}
$\dfrac{7}{11} < \dfrac{9}{11}$ (dezelfde \omterm{def:g4:fractions:def}{noemer}).

$\dfrac56 = \dfrac{15}{18} > \dfrac{11}{18}$.

$\dfrac{13}{12} > 1$ terwijl $\dfrac{19}{20} < 1$, dus
$\dfrac{13}{12} > \dfrac{19}{20}$.
\end{solution}

\begin{solution}{exo:g7:fractions:5}
$\dfrac59 + \dfrac29 = \dfrac79$; \quad
$\dfrac{11}{7} - \dfrac47 = \dfrac77 = 1$; \quad
$\dfrac58 + \dfrac78 = \dfrac{12}{8} = \dfrac32$.
\end{solution}

\begin{solution}{exo:g7:fractions:6}
$\dfrac23 + \dfrac56 = \dfrac46 + \dfrac56 = \dfrac96 = \dfrac32$.

$\dfrac{7}{10} - \dfrac25 = \dfrac{7}{10} - \dfrac{4}{10} = \dfrac{3}{10}$.

$\dfrac34 + \dfrac{5}{12} = \dfrac{9}{12} + \dfrac{5}{12} = \dfrac{14}{12} =
\dfrac76$.
\end{solution}

\begin{solution}{exo:g7:fractions:7}
$3 - \dfrac54 = \dfrac{12}{4} - \dfrac54 = \dfrac74$; \quad
$1 + \dfrac38 = \dfrac{11}{8}$; \quad
$2 - \dfrac76 = \dfrac{12}{6} - \dfrac76 = \dfrac56$.
\end{solution}

\begin{solution}{exo:g7:fractions:8}
$5 \times \dfrac{3}{10} = \dfrac{15}{10} = \dfrac32$; \quad
$\dfrac78 \times 4 = \dfrac{28}{8} = \dfrac72$; \quad
$\dfrac23$ van $45$ is $\dfrac{2 \times 45}{3} = \dfrac{90}{3} = 30$.
\end{solution}

\begin{solution}{exo:g7:fractions:9}
Samen: $\dfrac14 + \dfrac38 = \dfrac28 + \dfrac38 = \dfrac58$ van de pizza. Over:
$1 - \dfrac58 = \dfrac38$.
\end{solution}

\begin{solution}{exo:g7:fractions:10}
Uitgeschonken: $2 \times \dfrac16 = \dfrac26 = \dfrac13$ L. Over:
\[
\frac34 - \frac13 = \frac{9}{12} - \frac{4}{12} = \frac{5}{12}
\text{ L}.
\]
\end{solution}

\begin{solution}{exo:g7:fractions:11}
De positieve hoeveelheid $\frac12$ bij $\frac35$ optellen moet een uitkomst geven
die \emph{groter} is dan $\frac35$. Maar $\frac47 < \frac35$ (immers
$\frac47 = \frac{20}{35}$ en $\frac35 = \frac{21}{35}$): de beweerde \omterm{def:g1:addition:def}{som} is
kleiner dan een van haar termen --- onmogelijk.
\end{solution}

\begin{solution}{exo:g7:fractions:12}
Stap voor stap:
\[
\frac12 + \frac14 = \frac34, \qquad
\frac34 + \frac18 = \frac68 + \frac18 = \frac78, \qquad
\frac78 + \frac1{16} = \frac{14}{16} + \frac1{16} = \frac{15}{16}.
\]
De tussensommen $\frac12, \frac34, \frac78, \frac{15}{16}, \dots$ liggen elke keer
half zo ver van $1$: bij elke stap wordt het ontbrekende deel gehalveerd. De
sommen naderen $1$ zonder het ooit te bereiken.
\end{solution}

\begin{solution}{pb:g7:fractions:1}
\textbf{1.} (a) $\frac12 + \frac14 + \frac14 = \frac24 + \frac14 + \frac14 =
\frac44 = 1$: geldig. (b) $\frac14 + \frac18 + \frac18 + \frac12 = \frac28 +
\frac18 + \frac18 + \frac48 = \frac88 = 1$: geldig.

\textbf{2.} $\frac12 + \frac14 + \frac18 = \frac48 + \frac28 + \frac18 =
\frac78$: er ontbreekt één achtste noot.

\textbf{3.} Een maat is $1 = \frac{16}{16}$: zestien zestiende noten. Een achtste
is $\frac18 = \frac{2}{16}$: twee zestienden.

\textbf{4.} Gepunteerde halve: $\frac12 + \frac14 = \frac34$. Gepunteerde \omterm{def:g3:division:half}{kwart}:
$\frac14 + \frac18 = \frac28 + \frac18 = \frac38$.

\textbf{5.} De gepunteerde halve is $\frac34$ waard (vraag 4): precies één
walsmaat. Andere vullingen, bijvoorbeeld: \omterm{def:g3:division:half}{kwart} $+$ \omterm{def:g3:division:half}{kwart} $+$ \omterm{def:g3:division:half}{kwart}
($\frac14 + \frac14 + \frac14 = \frac34$), of \omterm{def:g3:division:half}{kwart} $+$ achtste $+$ achtste $+$
\omterm{def:g3:division:half}{kwart} ($\frac14 + \frac18 + \frac18 + \frac14 = \frac{2 + 1 + 1 + 2}{8} =
\frac68 = \frac34$).

\textbf{6.} $\frac38 + \frac18 + \frac14 = \frac38 + \frac18 + \frac28 =
\frac68 = \frac34$: er ontbreekt één kwartnoot ($\frac14$).

\textbf{7.} Gepunteerde achtste: $\frac18 + \frac{1}{16} = \frac{2}{16} +
\frac{1}{16} = \frac{3}{16}$. \omterm{def:g3:division:half}{Kwart}: $\frac{4}{16}$. De \omterm{def:g3:division:half}{kwart} duurt langer.

\textbf{8.} $\frac12 + \frac18 = \frac48 + \frac18 = \frac58$, dus
$\frac{10}{16}$: tien zestienden.

\textbf{9.} $1 - \frac{11}{16} = \frac{16}{16} - \frac{11}{16} = \frac{5}{16}$
van de maat is stilte.

\textbf{10.} Beide voorbeelden geven samen $\frac{2 + 1 + 1}{16} = \frac{4}{16}$:
geldig. Alle ritmes die vier zestienden vullen met stukken van $2$ en $1$:
\[
2{+}2, \quad 2{+}1{+}1, \quad 1{+}2{+}1, \quad 1{+}1{+}2, \quad
1{+}1{+}1{+}1 :
\]
vijf ritmes.

\textbf{11.} Een ritme dat $n$ zestienden vult, eindigt op een zestiende --- en
wat ervoor komt vult $n - 1$ --- of op een achtste --- en wat ervoor komt vult
$n - 2$. Elk ritme wordt zo precies één keer geteld, dus
$\text{aantal}(n) = \text{aantal}(n-1) + \text{aantal}(n-2)$. De tabel:
\begin{center}
\renewcommand{\arraystretch}{1.2}
\begin{tabular}{c|cccccccc}
$n$ & $1$ & $2$ & $3$ & $4$ & $5$ & $6$ & $7$ & $8$ \\
\hline
aantal & $1$ & $2$ & $3$ & $5$ & $8$ & $13$ & $21$ & $34$ \\
\end{tabular}
\end{center}
Een halve noot ($8$ zestienden) kan op $34$ manieren getrommeld worden. ($n = 4$
geeft de vijf ritmes van vraag 10 terug.)

\textbf{12.} Elk getal is de \omterm{def:g1:addition:def}{som} van de twee ervoor: $3 + 5 = 8$, $5 + 8 = 13$,
$8 + 13 = 21$, $13 + 21 = 34$ --- de regel die vraag 11 bewees, duizend jaar oud
en nu bekend als de regel van Fibonacci.

\textbf{13.} $\frac16 + \frac{1}{12} = \frac{2}{12} + \frac{1}{12} =
\frac{3}{12} = \frac14$: dezelfde duur als twee rechte achtsten
($\frac18 + \frac18 = \frac14$). De swingruil is geldig.

\textbf{14.} Elke noot van de triool duurt $\frac{1}{12}$: controle,
$3 \times \frac{1}{12} = \frac{3}{12} = \frac14$
(\cref{prop:g7:fractions:mult}). Voor een halve noot: elke noot duurt $\frac16$,
want $3 \times \frac16 = \frac36 = \frac12$.

\textbf{15.} Gepunteerde \omterm{def:g3:division:half}{kwart} $+$ \omterm{def:g3:division:half}{kwart} $+$ \omterm{def:g3:division:half}{kwart}: in achtsten is dat
$3 + 2 + 2 = 7$ achtsten $= \frac78$: geldig. Halve en kwartnoten zijn $4$ en $2$
achtsten waard --- beide \emph{\omterm{def:g3:numbers:evenodd}{even}} getallen. Elke \omterm{def:g1:addition:def}{som} van \omterm{def:g3:numbers:evenodd}{even} getallen is \omterm{def:g3:numbers:evenodd}{even},
maar een zevenachtstemaat heeft $7$ achtsten nodig, een \omterm{def:g3:numbers:evenodd}{oneven getal}: onmogelijk.
Om een \omterm{def:g3:numbers:evenodd}{oneven} maat te vullen, moet er iets \omterm{def:g3:numbers:evenodd}{oneven} verschijnen --- een achtste of
een gepunteerde noot.
\end{solution}
